\documentclass[11pt,a4paper]{article}
\usepackage[T1]{fontenc}
\usepackage[utf8]{inputenc}
\usepackage{lmodern}
\usepackage[spanish,es-nodecimaldot]{babel}
\usepackage[margin=2.2cm]{geometry}
\usepackage{array,booktabs,xcolor,listings}
\usepackage{hyperref}
\definecolor{azul}{HTML}{173B5E}
\definecolor{gris}{HTML}{F4F6F8}
\hypersetup{colorlinks=true,urlcolor=azul,linkcolor=azul}
\setlength{\parindent}{0pt}
\setlength{\parskip}{5pt}
\setlength{\emergencystretch}{3em}
\newcommand{\codigo}[1]{\texttt{\detokenize{#1}}}
\newcommand{\archivo}[1]{{\small\path{#1}}\par}
\lstset{language=Python,basicstyle=\ttfamily\footnotesize,
 backgroundcolor=\color{gris},frame=single,rulecolor=\color{black!15},
 keywordstyle=\color{azul}\bfseries,commentstyle=\color{black!55},
 stringstyle=\color{black!75},breaklines=true,columns=fullflexible,
 keepspaces=true,showstringspaces=false,tabsize=4,
 aboveskip=5pt,belowskip=5pt,
 literate={á}{{\'a}}1 {é}{{\'e}}1 {í}{{\'i}}1 {ó}{{\'o}}1 {ú}{{\'u}}1
 {ñ}{{\~n}}1 {ü}{{\"u}}1}

\hypersetup{pdftitle={Autenticación, permisos y tokens}}
\begin{document}
{\LARGE\bfseries\color{azul} Autenticación, permisos y tokens\par}
\smallskip
Los fragmentos conservan las llamadas del código; se omiten comentarios,
registros de diagnóstico y partes de las funciones.

\section*{Inicio de sesión y permisos}
\archivo{src/gui/app.py}
\begin{lstlisting}
self.email_msal_app = PublicClientApplication(
    client_id=MICROSOFT_CLIENT_ID,
    authority='https://login.microsoftonline.com/common',
    token_cache=self.email_token_cache.cache
)
\end{lstlisting}
MSAL identifica la aplicación mediante el Client ID. El inicio de sesión
ocurre en Microsoft, donde se aplican las políticas de acceso y MFA del
tenant. Este cliente de escritorio no utiliza un secreto de aplicación.

Los tres permisos son \textbf{delegados}: el acceso se realiza en nombre
del usuario autenticado.
\begin{center}
\small
\begin{tabular}{@{}p{3cm}p{3cm}p{9cm}@{}}
\toprule
Permiso & Servicio & Uso en la aplicación \\
\midrule
\codigo{User.Read} & Microsoft Graph & Consultar nombre, correo y alias del usuario. \\
\codigo{Mail.Send} & Microsoft Graph & Enviar mensajes y adjuntos desde su buzón. \\
\codigo{AllSites.Read} & SharePoint & Descargar el archivo configurado. Permite leer otros sitios a los que el usuario tenga acceso; no queda limitado a ese archivo. \\
\bottomrule
\end{tabular}
\end{center}

\subsection*{Autenticación interactiva}
\begin{lstlisting}
result = self.sharepoint_msal_app.acquire_token_interactive(
    scopes=scopes, prompt='select_account')
\end{lstlisting}
Esta llamada abre el navegador y MSAL incorpora PKCE, que vincula el
canje del código de autorización con la solicitud que lo originó.
\href{https://learn.microsoft.com/en-us/entra/msal/python/getting-started/acquiring-tokens#interactive-vs-silent}{Funcionamiento de MSAL}.

\subsection*{Retorno del correo en la ventana de configuración}
\begin{lstlisting}
redirect_uri = "http://localhost:62190/"
# ... apertura del navegador y recepción del código ...
result = app.acquire_token_by_authorization_code(
    auth_code[0], scopes=scopes, redirect_uri=redirect_uri)
\end{lstlisting}
Microsoft devuelve un código que se intercambia por tokens. La contraseña
no pasa por este manejador.

Este retorno manual no valida \codigo{state} ni configura PKCE.
Además, el servidor usa \codigo{TCPServer(("", 62190), ...)}, que escucha
en todas las interfaces. Para incorporar esas protecciones, debe usarse
el flujo interactivo de MSAL y un retorno limitado al equipo local.

\newpage
\section*{Perfil y envío de correo}
\archivo{src/gui/app.py}
\begin{lstlisting}
headers = {'Authorization': f'Bearer {access_token}'}
response = requests.get(
    'https://graph.microsoft.com/v1.0/me', headers=headers)
response.raise_for_status()
user_info = response.json()
email = user_info.get('mail') or user_info.get('userPrincipalName')
name = user_info.get('displayName')
\end{lstlisting}
\codigo{User.Read} permite identificar la cuenta conectada. Los alias se
consultan con \path{/me?$select=proxyAddresses}. Estas llamadas consultan
el perfil; no descargan mensajes del buzón.

\archivo{src/email_service/email_sender.py}
\begin{lstlisting}
MS_SCOPES = ['User.Read', 'Mail.Send']

headers = {
    'Authorization': f'Bearer {access_token}',
    'Content-Type': 'application/json'
}
# ... construcción de message ...
response = requests.post(
    'https://graph.microsoft.com/v1.0/me/sendMail',
    headers=headers,
    json=message
)
\end{lstlisting}
El token viaja en la cabecera de una petición HTTPS a Microsoft Graph.
El cuerpo contiene asunto, texto HTML, destinatarios, copias y adjunto.
\codigo{Mail.Send} autoriza el envío; no concede lectura de correos.
HTTP 202 significa que Microsoft aceptó la solicitud, no que el
destinatario ya la recibió.
\href{https://learn.microsoft.com/en-us/graph/api/user-sendmail}{Referencia de sendMail}.

\subsection*{Archivo adjunto y remitente}
\begin{lstlisting}
with open(attachment_path, 'rb') as f:
    file_data = f.read()
    file_name = os.path.basename(attachment_path)
    encoded_content = base64.b64encode(file_data).decode()
    attachment = {
        "@odata.type": "#microsoft.graph.fileAttachment",
        "Name": file_name,
        "ContentBytes": encoded_content
    }
message["message"]["attachments"] = [attachment]
\end{lstlisting}
Se lee el archivo local indicado y se incorpora al mensaje. Base64 es una
codificación, no cifrado; HTTPS protege el transporte hacia Microsoft.

El campo \codigo{from} puede indicar otra dirección, pero no otorga acceso
a ese buzón. Para enviar desde un buzón ajeno o compartido se requiere
adaptar los scopes a \codigo{Mail.Send.Shared} y disponer de la delegación
correspondiente en Exchange.
\href{https://learn.microsoft.com/en-us/graph/outlook-send-mail-from-other-user}{Permisos del remitente}.

\newpage
\section*{Lectura de archivos de SharePoint}
\archivo{src/gui/app.py}
\begin{lstlisting}
raw_site_url = self.config_manager.get_sharepoint_site_url()
from urllib.parse import urlparse
parsed_url = urlparse(raw_site_url)
base_domain = f"{parsed_url.scheme}://{parsed_url.netloc}"
scopes = [f"{base_domain}/.default"]
\end{lstlisting}
La dirección configurada determina el recurso para el que se pide el
token. \codigo{/.default} solicita los permisos delegados configurados
para SharePoint; en este registro se utiliza \codigo{AllSites.Read}.
El usuario también debe tener acceso al archivo.

\archivo{src/processing/file_processor.py}
\begin{lstlisting}
server_relative_path = quote(self.sharepoint_file_url)
endpoint = f"{self.site_url}/_api/web/GetFileByServerRelativeUrl('{server_relative_path}')/$value"
headers = {
    "Authorization": f"Bearer {access_token}",
    "Accept": "application/json;odata=nometadata"
}
response = requests.get(endpoint, headers=headers, stream=True)

if response.status_code == 200:
    with open(self.local_file_path, 'wb') as f:
        for chunk in response.iter_content(chunk_size=8192):
            f.write(chunk)
\end{lstlisting}
La petición GET descarga el contenido mediante la API REST de SharePoint.
Esta operación no modifica ni sube archivos. La copia se guarda en
\codigo{data/temp/} y puede reutilizarse en ejecuciones posteriores.

El token se envía a la URL configurada. Debe ser una dirección HTTPS del
SharePoint autorizado; este fragmento no valida una lista de dominios.
La copia descargada queda sujeta a los permisos del disco local.

\subsection*{Reutilización del token}
\archivo{src/email_service/email_sender.py}
\begin{lstlisting}
accounts = self.msal_app.get_accounts()
if accounts:
    result = self.msal_app.acquire_token_silent(
        self.MS_SCOPES, account=accounts[0])
    if result and 'access_token' in result:
        return result['access_token']
\end{lstlisting}
MSAL busca un token válido en la caché o intenta renovarlo. Si no obtiene
uno, el envío requiere iniciar sesión otra vez. La renovación sigue
sujeta a las políticas de Microsoft; cerrar el navegador no elimina los
tokens guardados.
\href{https://learn.microsoft.com/en-us/entra/msal/python/getting-started/acquiring-tokens#interactive-vs-silent}{Adquisición silenciosa de tokens}.

\newpage
\section*{Almacenamiento de tokens y cierre de sesión}
\archivo{src/gui/app.py}
\begin{lstlisting}
self.email_token_cache = SecureTokenCache(
    cache_name='email_token_cache')
self.sharepoint_token_cache = SecureTokenCache(
    cache_name='sharepoint_token_cache')
\end{lstlisting}
Correo y SharePoint mantienen cachés separadas. Cada una puede contener
tokens de acceso y de renovación, por lo que su contenido debe tratarse
como una credencial.

\archivo{src/utils/token_cache.py}
\begin{lstlisting}
serialized_cache = self.cache.serialize()
# ... división en fragmentos ...
keyring.set_password('my_app', fragment_key, fragment)
\end{lstlisting}
La caché se guarda mediante \codigo{keyring}, que puede usar el almacén
de credenciales de Windows. La protección depende del backend activo:
el código no lo fuerza ni lo comprueba. Dividir la caché en fragmentos
resuelve límites de tamaño, pero no añade cifrado.
\href{https://keyring.readthedocs.io/en/latest/}{Almacenamiento con keyring}.

\begin{lstlisting}
fragment = keyring.get_password('my_app', fragment_key)
# ... unión de los fragmentos ...
self.cache.deserialize(serialized_cache)
\end{lstlisting}
Al abrir la aplicación, los fragmentos se recuperan y reconstruyen en
memoria para que MSAL pueda reutilizar la sesión.

\subsection*{Cerrar sesión}
\archivo{src/gui/app.py}
\begin{lstlisting}
self.config_manager.clear_credentials()
self.email_token_cache.clear_cache()
self.sharepoint_token_cache.clear_cache()
\end{lstlisting}
Se eliminan los datos de sesión de la configuración y se solicita vaciar
las dos cachés. En \codigo{clear_cache()} se utilizan estas llamadas:
\begin{lstlisting}
self.cache.deserialize("")
# ... recorrido de los fragmentos almacenados ...
keyring.delete_password('my_app', fragment_key)
\end{lstlisting}
La limpieza afecta a la memoria y al almacén local. No revoca el
consentimiento en Entra ni cierra la sesión del navegador. Tampoco borra
archivos descargados o correos enviados. Si el almacén falla, la eliminación
puede quedar incompleta.

Los registros de diagnóstico pueden contener cuentas, destinatarios,
rutas y errores. El servidor del retorno manual también puede registrar
la URL con el código de autorización; esos registros no deben compartirse
sin depurarlos.
\end{document}
